% 04-walls.tex: the schema as the security boundary.
\section{The walls: the schema as the security boundary}
\label{sec:walls}

\motif{A promise enforced only by the people who run the system is not a promise.}

\analogy{The wall does not ask who is passing. It holds against the mayor, the clerk and the stranger alike, and it holds again after the town is rebuilt from its plans. A rule written into the wall does not depend on anyone remembering it.}

\subsection{Why the database is not plumbing}

In most systems the database stores what the application decided. In Polaris the database decides,
and the application is one client among several: the operator's web application, the command-line
tool, a bulk loader, a database shell, a restore from backup. A rule that lives in application code
binds only the clients that run that code. A rule enforced by a trigger, a check constraint or a
unique index binds every client, survives every restore, and cannot be bypassed by the next caller,
because the application role has no data-definition privilege with which to remove it. The
governing sentence of the architecture is \emph{by construction}: the system cannot be operated in a
way that violates the constraint, because the constraint is enforced before any application code
runs. \sImpl

The design principle behind it is older than Polaris: make invalid states impossible to represent
rather than asking developers not to create them. A verification event that records a token
identifier at the zero-knowledge disclosure level is not an application bug to be caught in review.
It is a privacy violation that must never be storable, so it is a check constraint on the row.

\subsection{The ten constraints and where each lives}

The constitution names ten hard constraints in two tiers. Five are constitutional, rights
guarantees that may not be relaxed without an amendment recorded in the changelog; five are
engineering invariants, the disciplines that keep the first five true under load and attack. The
hierarchy is one line deep, vocation, then constitutional constraints, then engineering invariants,
and the constitution states that no deeper apparatus governs them.

\begin{table}[H]
\centering\footnotesize
\caption{The ten constraints at \code{1.0.0-rc.7}, their tier, and the object that enforces each. The table is kept honest against the code by \code{check\_c1c10\_objects\_resolve}, which fails the build if a named object no longer exists.}
\label{tab:constraints}
\setlength{\tabcolsep}{4pt}\renewcommand{\arraystretch}{1.1}
\rowcolors{2}{tabstripe}{white}
\begin{tabularx}{\textwidth}{>{\bfseries\hspace{0pt}}p{0.3in} >{\RaggedRight\arraybackslash\hspace{0pt}}X >{\RaggedRight\arraybackslash\hspace{0pt}}p{0.95in} >{\RaggedRight\arraybackslash\hspace{0pt}}p{2.0in}}
\toprule
\rowcolor{tabheader}\thead{\#} & \thead{Constraint} & \thead{Tier} & \thead{Enforced by} \\
\midrule
C1 & The audit trail is append-only; history is never rewritten or deleted & Constitutional & Triggers reject UPDATE and DELETE on every audit table (\code{reject\_audit\_modification}) \\
C2 & A zero-knowledge verification stores no token identifier & Constitutional & Bidirectional CHECK \code{chk\_disclosure\_token\_consistency} \\
C3 & One person holds at most one ACTIVE token & Constitutional & Partial unique index \code{uq\_one\_active\_per\_person} \\
C4 & Failed-login counting is atomic & Engineering & Single-statement \code{UPDATE \ldots RETURNING} \\
C5 & No inline scripts; the content security policy is \code{script-src 'self'} & Engineering & Response header, verified per route \\
C6 & Disclosure level is enforced server-side; a client cannot upgrade what it learns & Constitutional & Server code paired with redaction tests; the C2 CHECK behind it \\
C7 & No hardcoded cryptography; algorithms are rows in a registry & Engineering & Foreign key to \entity{CryptographicAlgorithm} \\
C8 & Every map and API aggregate is bounded & Engineering & Hard caps in the SQL functions and routes \\
C9 & Concurrency claims are tested with real threads & Engineering & Threaded suites against a live database \\
C10 & Identity is not money; the schema carries no monetary claim & Constitutional & Structural absence, pinned by a check \\
\bottomrule
\end{tabularx}
\end{table}

The repository also records a structural claim about the ten: that the set is closed, so that an
eleventh constraint requires either replacing one or extending the framework, and that the
constraints depend on one another, so that removing any one ripples through the rest. The
dependency walk (remove C10 and every other constraint protects a financial-surveillance system
instead; remove C1 and every other constraint's enforcement becomes retroactively deniable) is the
argument for treating them as a graph rather than a list. Athena, in Section~\ref{sec:cognition},
holds the map from each constraint to the exact live mechanism that enforces it, and a check fails
the build if a mechanism drifts. Three of the ten now also have a formal model checked on every push:
C1's purge coverage, C2's unlinkability across two pooling relying parties, and C3 under concurrent
issue and revoke, each with a counterpart configuration it must fail, so none of the three is vacuous.

\figpage{districts}{The 45 canonical tables of \code{01\_schema.sql} at \code{1.0.0-rc.7}, as rings
around the nucleus. The faint lattice at the centre is the nucleus geometry of
Figure~\ref{fig:nucleus} at the same proportions. Ring two is the credential machinery: signatures,
permissions, epochs, recovery, duress, the holder's key, the card. Ring three is authority, trust and
the people at the counter: who may sign under which algorithm, who trusts whom, revocation ceilings,
quotas, relying parties, the replay registers, and what an enrolment rested on. Ring four is the
archives, watchtowers and decisions kept as records: append-only logs, anchors, the retention decision
as data, and the events that record every change to an authority, an operator account or a relying
party. A migrated deployment holds seven more tables, drawn dashed: the migration registry closes the fourth ring, so rings two, three and four have the same thirteen spokes, and the other six (operator sessions, their hardware keys, the audit-access log and Athena's three curated tables) form a fifth, hexagonal ring that echoes the nucleus. Its arrows all turn one way: migrations applied one after another, round the canonical schema. The five partition children are not drawn.}{fig:districts}

\subsection{The audit of record}

Fourteen tables follow one named pattern, the audit of record: the row and its history are the same
object. There is no separate event log beside the table that could be edited while the table is
locked. Mutation is either forbidden outright or bounded to one field that moves one way, such as a
signature's deprecation date or an attestation's revocation date. Every one of the fourteen is
enforced by a trigger that raises \code{insufficient\_privilege}, and a check fails the build if any
trigger stops doing so. Beyond the fourteen the design record names, thirty-one tables in all carry
an append-only trigger at \code{1.0.0-rc.7}, the new ones being the records that decisions about
authorities, operators, relying parties and enrolments now leave behind, and C1 is
tested once per table that claims it. Of the thirty-one, twenty-four refuse every edit, and a
table-driven suite tests each of them for both UPDATE and DELETE, cross-checked against the
catalogue; the other seven permit bounded mutation, one field moving one way, and the trigger
mutation drill is what answers for them. The cross-check finds tables by their guard function, and
the list of those functions was written by hand, so a new table under a new guard would have gone
untested without anything saying so. Since 23 September 2026 every such guard must be declared
strict or bounded, which is deciding which instrument answers for it, and an undeclared one fails
the build. \sImpl

The corollary is that no foreign key anywhere in the schema cascades a delete. A cascade would let
the lifecycle events of a deleted token vanish with it. Every reference defaults to \code{NO
ACTION}, so a principal referenced by any audit row is effectively undeletable, and a check scans
every SQL file for both cascade forms with no allowlist. The one record that deliberately carries no
foreign key at all, so that ending a relying party's contract cannot erase its history, taught the
repository something about its own test database at \code{1.0.0-rc.7}: the seed that reloads the
sample world restarted the party's identifiers from one and reached that record with neither, so
the first party of each reload inherited a dead party's events. The seed names the record now, and a
check computes from the schema's own foreign keys which records a reload would leave behind.

The application role, \code{polaris\_app}, is revoked UPDATE and DELETE on the append-only tables
and holds no data-definition privilege. The only path that deletes audit rows is the archive-then-purge
procedure, which verifies the archive against its manifest first, refuses any cutoff inside the
retention window, and appends its own append-only checkpoint row. The retention window itself is
data with a floor: no policy row can express a retention shorter than a year, because a check
constraint refuses it, and shortening the floor is a schema change answered for in review. The
vocation refuses unbounded retention; it equally refuses retention short enough to erase the record.
\sImpl

\subsection{The check layer, and the drills that attack it}

Above the schema sits a flat layer of 313 plain functions, each of the form
\code{check\_*(repo\_root)}, that read the repository and return a finding. They gate continuous
integration: any failure fails the build. Each check is paired with a detection test that breaks a
fixture in the specific way the check claims to catch and asserts that the check turns red, and also asserts that the check passes on the same fixture unbroken, because a test
that only ever asserts failure has not established that it detects anything. A check that cannot
detect its own violation is treated as broken. This is the mechanism by which prose claims are turned
into gates: a count stated in the README, a route that must return a header, a migration that must
not cascade, a drill that CI must run. \sImpl

The walls are tested as well as trusted. A mutation drill removes each check, each CHECK constraint,
each trigger, each unique index and each refusal a stored procedure makes, one at a time, and requires
something to go red; it runs on every push, with a negative control so that zero survivors cannot be
confused with a harness that ran nothing. Section~\ref{sec:senses} describes that layer, and
\evidencerecord{} records what it found, including the suite that exists to prove C3
violating C3.
\sImpl

\begin{layercard}{Layer card: the walls}
\qa{What it is}{PostgreSQL 16: 45 tables, 46 triggers (42 of them in the trigger file), sixteen stored procedures and four use-case functions that implement every state-changing use case, the append-only audit tables, the retention policy as data.}
\qa{Why it exists}{A promise enforced only by the people who run the system is not a promise. The schema binds every client, every restore and every insider with a database connection short of the owner.}
\qa{What it trusts}{The database engine and its owner; the correctness of the constraints as written, which is why each is now mutation-tested. Above the owner there is nothing: an owner can disable a trigger, drop a policy or rewrite a row, and every guarantee in this section falls with it. Against the owner the only defence is evidence held outside the database, the transparency logs and their witnesses, and no institutionally independent witness exists (Section~\ref{sec:limitations}); until one does, the schema binds every actor except its sovereign. It also binds only the roles that actually run: several of the corrections recorded in \evidencerecord{} since this map was drawn are guarantees that held for the owner, whom the suites connected as, and not for the application role production uses.}
\qa{What it refuses}{Application code, review process and operator discipline as the last line: the application never enforces a constitutional guarantee on its own. A test that has never been shown to fail.}
\qa{What it can see}{Everything the system stores. This is why the privacy layer bounds what is stored, not who may read it.}
\qa{What it cannot see}{Biometric templates, holder-side keys, enrolment documents, the content of exchanged messages or timestamped documents: none enter the database.}
\qa{If it fails}{A dropped trigger fails the build: deleting each of the 42 in the trigger file in turn, a check notices 38 by reading the schema, and the trigger mutation drill notices all 42 by running the suites; a lowered retention floor or a cascade fails the checks; a constraint whose loss the suite would not notice fails the mutation drill; a database owner with a shell has larger options than any trigger, which the threat model states rather than hides.}
\qa{Now or later}{Now.}
\qa{Inside and outside}{Inside: the credential core. Outside: cryptographic evidence, because a constraint binds only the database that holds it.}
\end{layercard}


\nextlayer{A constraint binds the database, not a copy of the claim that has left it. The moment a
credential is handed to a stranger, or an authority's history is read by an auditor, nothing in the
schema travels with it. The next layer makes claims carry their own proof.}
